\subsection{Principal Bundles}

\begin{definition}
    \label{fiber_bundles:def:principal_bundle}
    Let $G$ be a topological group. A $G$-principal bundle is a fiber bundle with fiber $G$ and structure group $G$ acting by left translations.
\end{definition}

The total space of a $G$-principal bundle $P \xrightarrow{\pi} B$ carries a natural right action of $G$ defined as follows. Around $y \in P$ choose a local trivialisation
\[
\phi \colon U \times G \to \pi^{-1}(U)
\]
and let $\phi^{-1}(y) = (\pi(y),h)$. Then set 
\[
y \cdot g := \phi(\pi(y),hg).
\]
This is well defned since by definition $G$ acts on itself by \textit{left} translations under a change of local trivialisation. We then have the following Lemma.

\begin{lemma}
    Let $P \xrightarrow{\pi} B$ be a $G$-principal bundle. Then $\pi$ descends to a homeomorphism $\modulo{P}{G} \to B$.
\end{lemma}

\begin{proof}
    Clearly $\pi$ descends to a continuous bijection 
    \[
    \hat{\pi} \colon \modulo{P}{G} \to B.
    \]
    To show that $\hat{\pi}$ is in fact a homeomorphism let us first assume that $P$ is trivial. Let $\tau\colon B \times G \to P$ be a trivialisation. Consider the commutative diagram
   \[
\begin{tikzcd}[every label/.append style={font=\normalsize}]
	& {B \times G} & P & B \\
	B & {\modulo{(B \times G)}{G}} & {\modulo{P}{G}}
	\arrow["\tau", "{\scriptstyle\cong}"', from=1-2, to=1-3]
	\arrow["p"', from=1-2, to=2-1]
	\arrow[from=1-2, to=2-2]
	\arrow["\pi", from=1-3, to=1-4]
	\arrow[from=1-3, to=2-3]
	\arrow["\hat{p}"', "{\scriptstyle\cong}", from=2-2, to=2-1, pos = 0.3]
	\arrow["{\hat{\tau}}", "{\scriptstyle\cong}"', from=2-2, to=2-3]
	\arrow["{\hat{\pi}}"', from=2-3, to=1-4]
\end{tikzcd}
\]
where the action of $G$ on $B \times G$ corresponds to the action of $G$ on $P$ under $\tau$ and is explicitly given by 
\[
(x,h) \cdot g = (x,hg).
\]
Moreover $\hat{\tau}$ is the map induced on the quotients by $\tau$ and $\hat{p}$ is the homeomorphism induced by the projection $p \colon B \times G \to B$. Then for the lower route in the diagram we have 
\[
\hat{\pi} \circ \hat{\tau} \circ \hat{p}^{-1} = \mathrm{id}_B.
\]
Since $\hat{\tau}$ and $\hat{p}^{-1}$ are homeomorphisms, so is $\hat{\pi}$. Now for the general case note that the local triviality of $P$ together with the above shows that $\hat{\pi}$ is a local homeomorphism. Since it is bijective the statement follows.
\end{proof}

A morphism of principal bundles should respect this right action, thus we define:

\begin{definition}
    A morphism from a $G$-principal bundle $P \to B$ to an $H$-principal bundle $P^\prime \to B^\prime$ is a morphism $(F,f) \colon (P \to B) \to (P^\prime \to B^\prime)$ of the respective fiber bundles (see Definition \ref{fiber_bundles:def:bundle_morphism}) together with a homomorphism $\phi \colon G \to H$ of topological groups such that $F$ is $\phi$-equivariant, that is $$F(x \cdot g) = F(x) \cdot \phi(g)$$ for all $x \in P$ and $g \in G$. If $G = H$ we usually assume $\phi$ to be the identity and if $B = B^\prime$ we assume $f$ to be the identity.
\end{definition}


The added condition of equivariance is actually quite strong as illustrated by the following proposition.

\begin{proposition}
    \label{fiber_bundles:prop:morphisms_of_principal_bundles}
    Any morphism of $G$-principal bundles over the same base space is an isomorphism.
\end{proposition}

\begin{proof}
    First suppose we are given a morphism 
    \[  
        F \colon B \times G \to B \times G
    \]
    of trivial bundles. Then by definition $F$ is of the form 
    \[
    F(x,g) = (x,\phi_x(g))
    \]
    for a continuous family of equivariant maps $\phi_x \colon G \to G$. By equivariance $\phi_x(g) = \phi_x(e) g$ and thus an inverse of $F$ is given by
    \[
    F^{-1}(x,g) := (x, \phi_x(e)^{-1} g).
    \]
    Since every $G$-principal bundle is locally trivial this shows that we can always find local inverses in the general case and since inverses are unique they fit together to define a global inverse.
\end{proof}

\begin{corollary}
    A $G$-principal bundle is trivial if and only if it admits a section.
\end{corollary}

\begin{proof}
    Suppose we are given a section $s \colon B \to P$ of $P \to B$. Then $B \times G \to P, \; (x,g) \mapsto s(x) \cdot g$ defines a morphism of $G$-principal bundles and hence is an isomorphism by Proposition \ref{fiber_bundles:prop:morphisms_of_principal_bundles}.
\end{proof}

\needspace{5cm}

\begin{remark}
    In the defintion of $G$-principal bundles we may also require $G$ to act by right translations instead of left translations on itself. Correspondingly we then get an induced left $G$-action on the total space. Let us call those bundles right $G$-principal bundles and the bundles from Defintion \ref{fiber_bundles:def:principal_bundle} left $G$-principal bundles. Given a right $G$-principal bundle $P \to B$ we can always turn it into a left $G$-principal bundle as follows. Let $\tau \colon U \times G \to P$ be any local trivialisation and set 
    \[
    \tau^\prime := U \times G \xrightarrow{\mathrm{id} \times \mathrm{inv}} U \times G \xrightarrow{\tau} P,
    \]
    where $\mathrm{inv} \colon G \to G$ is the inversion map. Then the maps $\tau^\prime$ constructed in this way give $P \to B$ the structure of a left $G$-principal bundle. The new right $G$-action on $P$ can be expressed by the left $G$-action as 
    \[
    y \cdot g = g^{-1} \cdot y
    \]
    Note that any morphism of right principal bundles also respects this new right action and thus defines a morphism of the corresponding left prinicpal bundles. In other words, this construction defines an isomorphism between the categroy of right principal bundles and the category of left principal bundles. Thus there is virtually no difference between these two notions.
\end{remark}

Given any fiber bundle we have seen in Corollary \ref{fiber_bundles:cor:changing_fiber}, that we can change the fiber. If we start with a principal bundle we can make this construction quite explicit:

\begin{lemma}[Borel Construction]
    \label{fiber_bundles:lem:borel_construction}
    Let $P \xrightarrow{p} B$ be a $G$-principal bundle and $F$ a space upon which $G$ acts effectively from the left. Then 
    \[
    P \times_G F := \modulo{(P \times F)}{\sim}, \quad (y,f) \sim (y\cdot g, g^{-1}f)
    \]
    together with the map $\pi \colon P \times_G F \to B$ given by $[(y,f)] \mapsto p(y)$ is a fiber bundle with fiber $F$ and structure group $G$. Moreover this construction yields the same bundle as Corollary \ref{fiber_bundles:cor:changing_fiber}, hence the same notation.
\end{lemma}

\begin{proof}
    The map $\pi$ is induced by $P \times F \to B, \; (y,f) \mapsto p(y)$ and hence continuous. To see that $P \times_G F$ is locally trivial let $\lambda \colon U \times G \to P$ be a local trivialisation. Then the homeomorphism 
    \[
    \lambda \times \mathrm{id}_F \colon U \times G \times F \to p^{-1}(U) \times F
    \]
    descends to a homeomorphism
    \[
    \tilde{\kappa} \colon \modulo{(U \times G \times F)}{\sim} \to \modulo{(p^{-1}(U) \times F)}{\sim} = \pi^{-1}(U),
    \]
    where the equivalence relation on $U \times G \times F$ is given by $(x,g,f) \sim (x,gh,h^{-1}f)$. Moreover the map 
    \[
    U \times G \times F \to U \times F, \; (x,g,f) \mapsto (x,gf)
    \]
    descends to a homeomorphism $\modulo{(U \times G \times F)}{\sim} \to U \times F$ with inverse given by $(x,f) \mapsto [(x,e,f)]$. Then under this homeomorphism $\tilde{\kappa}$ corresponds to a homeomorphism $\kappa \colon U \times F \to \pi^{-1}(U)$, which is explicitly given by $(x,f) \mapsto [(\lambda(x,e),f)]$ and
    \[\begin{tikzcd}
	{U \times F} & {\pi^{-1}(U)} \\
	{} & U
	\arrow["{\kappa}", from=1-1, to=1-2]
	\arrow["{\mathrm{proj.}}"', from=1-1, to=2-2]
    \arrow["\pi", from=1-2, to=2-2]
    \end{tikzcd}\]
    commutes. So $P \times_G F$ is locally trivial. Next we have to compute the transition functions. Let $U_\alpha, U_\beta \subseteq B$ be two trivialising neighbourhoods for $P$ with  local trivialisations $\lambda_\alpha, \lambda_\beta$ and transition function $\tau_{\alpha,\beta} \colon U_\alpha \cap U_\beta \to G$. Moreover let $\kappa_\alpha, \kappa_\beta$ be the corresponding local trivialisations for $P \times_G F$ as constructed above. Then we compute 
    \begin{align*}
        (\kappa_\alpha^{-1} \circ \kappa_\beta)(x,f) & = \kappa_\alpha^{-1}([(\lambda_\beta(x,e),f)]) = (x,\tau_{\alpha,\beta}(x) f).
    \end{align*}
    So $P \times_G F$ has structure group $G$ and from the uniqueness result of Theorem \ref{thm:fiber_bundle_construction1} it follows that this construction agrees with the one from Corollary \ref{fiber_bundles:cor:changing_fiber}.
\end{proof}

Given a group homomorphism $\phi \colon G \to H$ of topological groups we can view $H$ as a left $G$-space by $g \cdot h := \phi(g)h$. Of course this action may not be effective anymore but the Borel construction can still be used to construct an $H$-principal bundle out of a $G$-principal bundles using this action:

\begin{proposition}
    \label{fiber_bundles:prop:borel_construction_part2}
    Let $P \xrightarrow{p} B$ be a $G$-principal bundle and $\phi \colon G \to H$ a homomorphism of topological groups. Then 
    \[
    P \times_\phi H := \modulo{(P \times H)}{\sim}, \quad (y,h) \sim (y\cdot g, \phi(g)^{-1}h)
    \]
    together with the map $\pi \colon P \times_\phi H \to B$ given by $[(y,h)] \mapsto p(y)$ is an $H$-principal bundle. Moreover this construction is natural in the sense that if $f \colon B^\prime \to B$ is a continuous map then 
    \[
    (f^\ast P) \times_\phi H \cong f^\ast (P \times_\phi H)
    \]
    and it is also functorial in the sense that if $\psi \colon H \to K$ is another homomorphism of topological groups then 
    \[
    (P \times_\phi H) \times_\psi K \cong P \times_{\psi \circ \phi} K.
    \]
\end{proposition}

\begin{proof}
    The proof is the same as the one for Lemma \ref{fiber_bundles:lem:borel_construction}, for the special case $F = H$ and the specific left action of $G$ on $H$ determined by $\phi$. We only used the effectiveness of the group action at the very end to make use of Theorem \ref{thm:fiber_bundle_construction1}, which only makes sense if the group action is effective, such that the transition functions are uniquely determined (and besides that we required the group action to be effective in the definition of a fiber bundle, so the statement of the Lemma wouldn't even technically make sense). The difference in the statement of this proposition is that $P \times_\phi H$ will be a bundle with structure group $H$ and not $G$ and $H$ does again act effectively on itself. Using the notation of Lemma \ref{fiber_bundles:lem:borel_construction} we then see that changing local trivialisations is given by 
    \[
        (\kappa_\alpha^{-1} \circ \kappa_\beta)(x,h) = (x,\phi(\tau_{\alpha,\beta}(x)) h),
    \]
    such that $P \times_\phi H$ is indeed an $H$-principal bundle. The naturality and functoriality statements are shown analogously to Proposition \ref{fiber_bundles:prop:natruality_of_changing_fiber} and we shall use the same terminology of a trivialising cover introduced in that proof. Then the construction of the local trivialisations for $P \times_\phi H$ shows that any trivialising cover $(U_\alpha)$ for $P$ with transition functions $\tau_{\alpha,\beta}$ is also a trivialising cover for $P \times_\phi H$ with transition functions 
    \[
    \phi \circ \tau_{\alpha,\beta} \colon U_\alpha \cap U_\beta \to H.
    \]
    Hence we see as in Proposition \ref{fiber_bundles:prop:natruality_of_changing_fiber} that both $(f^\ast P) \times_\phi H$ and $f^\ast (P \times_\phi H)$ have $(f^{-1}(U_\alpha))$ as a trivialising cover with transition functions 
    \[
    \phi \circ \tau_{\alpha,\beta} \circ f \colon f^{-1}(U_\alpha) \cap f^{-1}(U_\beta) \to H,
    \]
    so the naturality result follows again from Theorem \ref{thm:fiber_bundle_construction1}. For the functoriality note that both bundles have $(U_\alpha)$ as a trivialising cover with transition functions $\psi \circ \phi \circ \tau_{\alpha,\beta}$.
\end{proof}

If $\mathcal{P}_G(B)$ denotes the set of isomorphism classes of $G$-principal bundles over a given base space $B$ then $\mathcal{P}_G(-)$ defines a contravariant functor $\mathsf{Top} \to \mathsf{Set}$. Then the above proposition says that any $\phi \colon G \to H$ defines a natural transformation $\mathcal{P}\phi := \times_\phi H$ of functors $\mathcal{P}_G(-) \Rightarrow \mathcal{P}_H(-)$ and does so in a functorial way, that is $\mathcal{P}\psi \circ \mathcal{P}\phi = \mathcal{P}(\psi \circ \phi)$. Thus we get a functor from the category of topological groups into the functor category of functors $\mathsf{Top} \to \mathsf{Set}$.
